% SPDX-License-Identifier: LPPL-1.3c
% Copyright (C) 2026 Christophe Jorssen
% Author and Current Maintainer: Christophe Jorssen <christophe.jorssen@gmail.com>
% LPPL maintenance status: maintained
% This file is part of luafemm. See LICENSE and LICENSES.md for its terms.
\section{A rotor moves inside a circular air band}
\label{sec:sliding}

\subsection{Kamil turns the rotor}
Kamil has drawn his machine. He now wants to turn its rotor while keeping the
stator current fixed, plot the magnetic field at a few positions, and compare
torque with flux linkage. Redrawing and remeshing the entire machine at every
angle would mix physical changes with changes in the finite-element mesh.
A \emph{moving air band} keeps the bulk triangles attached to their respective
bodies. Only a narrow annulus of air receives new triangles at each position.

This first implementation treats independent \emph{planar static} positions.
The complete nonlinear field is solved afresh at each angle. It introduces
neither angular velocity nor motion-induced voltage. Its triangles are ordinary
P1 finite elements, not FEMM's specialised analytical air-gap elements (AGE).
There can still be discretisation ripple when the band diagonals change.

\Needspace{10\baselineskip}
\subsection{Declare once, capture, then choose positions}
The order matters. First declare the problem, materials, circuits and all
physical paths or nodes. Then capture a mesh-only template. Finally create
one new problem per requested angle and perform the measurements there.
\begin{enumerate}
\item Draw every rotor domain strictly inside an inner circle. Draw every
  stator domain outside an outer circle. Leave a source-free air collar on
  both sides of the circles: neither trace may coincide with a material edge.
\item Attach |femm/sliding band| to the declaration picture. The key captures
  at the end of that picture, after all its paths. Do not request a field,
  mesh, profile value, measurement or operating point before this capture.
\item In a new picture, set |femm/problem={sliding band=machine,rotor angle=30}|.
  This copies the physical data, turns the rotor and connects the band.
\item Draw the inherited geometry/field and measure torque, circuit linkage
  or ordinary profiles. Repeat step 3 for another angle.
\end{enumerate}

\begin{codeexample}[code only]
\tikzset{machine position/.style={
  femm/problem={sliding band=machine,rotor angle=#1},
  femm/cache={file=machine-angle-#1,mode=auto}}}
% A preceding declaration picture captured the template named machine.
\begin{tikzpicture}[machine position=30]
  \pic {femm geometry};
  \pic {femm field};
  \femmmeasure[name=at-30,quantity=force,method=band,parameter=30]
  \femmmeasure[name=port-30,quantity=circuit,circuit=drive,parameter=30]
  \path[femm/profile={name=gap-30,samples=181}]
    (4.8,0) arc[start angle=0,end angle=360,radius=4.8];
\end{tikzpicture}
\end{codeexample}
The picture transformation controls presentation. It does not turn the physical
rotor: that is the job of |rotor angle|. In particular, an untagged rotated
ellipse is just a drawing overlay. Tagged physical paths in the declaration
picture retain their full TikZ expressiveness; the usual flattening tolerance
still controls their computational polygons. The band traces themselves are
regular polygonal approximations of circles.

\Needspace{10\baselineskip}
\subsection{A complete reluctance-machine experiment}
Kamil uses an elliptical rotor and an annular stator with two prescribed
100-turn winding legs. The signs of their turns close the winding, so its
linkage is independent of the potential gauge. The steel table is illustrative
and single-valued. The two field panels use independently chosen contour levels;
compare the measured linkage, rather than visible line counts, for magnitude.
The dashed circles show the artificial band, not a new material interface.

The source starts with a relatively coarse mesh. While adjusting colours and
labels, Kamil uses only two or three angles and |mesh size=1.2|. He then decreases
the bulk spacing, increases |segments| and checks several |layers| values.
For a final quantitative plot he also moves the exterior box farther away.
A smooth-looking interpolated plot is not evidence of converged torque.

\luafemmexample{sliding-band}

At fixed current, the rotor tends towards increasing magnetic coenergy. For
this example it aligns approximately at $90^\circ$ and the torque reverses
there. The geometry repeats after $180^\circ$; the API itself normalises angles
only modulo $360^\circ$. Asymmetry in a discretised template can cause small
differences between nominally equivalent half-turn positions.

The loop in this example is ordinary |\foreach|. Each named result remains
available to |\femmresultplot| after its picture has ended. Store the intended
abscissa explicitly with |parameter|. Changing |femm/lines|, colour, scale or
profile sample count does not change the solved physical position.

\Needspace{10\baselineskip}
\subsection{Capture keys and geometric contract}
\Needspace{7\baselineskip}
\begin{key}{/tikz/femm/sliding band=\marg{key list}}
Schedule capture at the end of this picture, binding it to this picture's model.
The keys belong to |/femm/sliding band|. There is no default template name.
Only one band is supported per template. All physical declarations must precede
the capture and no mesh or solve may already exist. A mesh cache restored
\emph{by the capture itself} is allowed.
\end{key}
\Needspace{7\baselineskip}
\begin{command}{\femmslidingband\oarg{key list}}
Capture immediately. Place this command after the last physical declaration.
It accepts the same key family as the picture key. It does not draw or solve.
Each invocation resets the capture options to the defaults below before applying
its list; package styles can supply recurring settings.
\end{command}
\Needspace{7\baselineskip}
\begin{key}{/femm/sliding band/name=\meta{name} (initially empty)}
Required document-wide template identifier. Duplicate names are errors. The
name remains available across TeX groups for the current run; it is not itself
a saved file or a persistent object identifier.
\end{key}
\Needspace{7\baselineskip}
\begin{key}{/femm/sliding band/model=\meta{name} (initially empty)}
For the immediate command, select an existing named model instead of the current
one. It must still be unmeshed. The picture-level key supplies its own originating
model name after the user's list, so later pictures cannot redirect its capture.
\end{key}
\Needspace{7\baselineskip}
\begin{key}{/femm/sliding band/inner radius=\meta{number} (initially 0)}
Inner trace radius in model units. Required positive value. Every rotating path
must lie strictly inside the inscribed polygon, not merely inside the circle.
\end{key}
\Needspace{7\baselineskip}
\begin{key}{/femm/sliding band/outer radius=\meta{number} (initially 0)}
Outer trace radius in model units, strictly larger than the inner radius.
Every stator path must lie outside this circle; compound stator paths may
surround it but their material must not fill the band. The entire circle must
lie strictly inside the rectangular problem box.
\end{key}
\Needspace{7\baselineskip}
\begin{key}{/femm/sliding band/center x=\meta{number} (initially 0)}
Horizontal coordinate of the rotation centre in model units.
\end{key}
\Needspace{7\baselineskip}
\begin{key}{/femm/sliding band/center y=\meta{number} (initially 0)}
Vertical coordinate of the same centre. These are model coordinates, unaffected
by later drawing transformations. They do not set the torque origin.
\end{key}
\Needspace{7\baselineskip}
\begin{key}{/femm/sliding band/segments=\meta{integer} (initially 96)}
Number of vertices on each initially declared trace and each intermediate ring;
allowed range 16--1024. The template mesher may subdivide the trace chords,
so their final vertex counts can differ. Raising this value refines the angular
resolution and the polygonal geometry; it does not refine the bulk interiors.
\end{key}
\Needspace{7\baselineskip}
\begin{key}{/femm/sliding band/layers=\meta{integer} (initially 2)}
Number of radial strips, from 1 to 16. Intermediate rings are equally spaced
in radius and fixed in the stator frame. With inner radius $r_i$, outer radius
$r_o$, $L$ layers and $N$ segments, the nesting check requires
\[
 \bigl(r_i+(r_o-r_i)/L\bigr)\cos(\pi/N)>r_i.
\]
If it fails, increase |segments| or reduce |layers|. Excessive radial refinement
with insufficient angular refinement creates poor triangles. The ordinary
minimum-angle and maximum-area targets, when nonzero, are checked on every
new position; an unmet target is an error, not a silently accepted mesh.
\end{key}
All six numerical capture keys accept ordinary PGF mathematical expressions.
No Lua source is evaluated from key values. A path crossing or filling the band
is rejected, including a large air refinement rectangle. Use local spacing
on rotor/stator regions instead. Both traces must be in air with vacuum
permeability, zero coercivity and no current. The region order retains its
usual meaning inside each bulk domain.

\Needspace{10\baselineskip}
\subsection{Position and measurement keys}
\Needspace{7\baselineskip}
\begin{key}{/femm/problem/sliding band=\meta{name} (initially empty)}
Create an independent position from the named template. An empty value selects
the ordinary problem workflow. The template owns geometry, materials, circuits,
exterior conditions, mesh targets, units, depth and solver settings. Do not
redeclare them in a position picture. Graphics, profiles, selections and
measurements may still be added. Cache options belong to the new picture.
\end{key}
\Needspace{7\baselineskip}
\begin{key}{/femm/problem/rotor angle=\meta{number} (initially 0)}
Mechanical rotation in degrees, counterclockwise about the band centre, relative
to the captured geometry. PGF expressions are allowed, for example |90/3|.
Negative values are allowed and finite angles are reduced modulo $360^\circ$.
A nonzero angle without |sliding band| is an error. It is not an electrical
angle or a display phase. Ordinary TeX grouping controls problem-key inheritance.
\end{key}
The existing |/femm/problem/circuit currents| list can override prescribed
series currents in amperes. Unlisted currents retain their template values.
This differs from the zero-excitation default of an incremental response.
Unknown circuit names and phase overrides are rejected. The new current and
angle enter the physical cache signature. No other physical problem override
is accepted, even if it would numerically equal the inherited setting.

\Needspace{7\baselineskip}
\begin{key}{/femm/measure/method=band}
Additional choice for |quantity=force| on a sliding position. Integrate Maxwell
stress across the entire air band to obtain the total |Fx|, |Fy| and |torque|
on everything inside it. No body selection or profile is needed or accepted.
The origin remains |origin x=0,origin y=0| in model units; explicitly set it to
the band centre for torque about an offset rotation axis. Forces use newtons,
torque uses newton metres, and positive torque is counterclockwise. The problem
must have positive |depth|. Existing integral components cover the whole model.
\end{key}
The general measurement default remains |method=weighted|. Set |method=band|
explicitly, or define a style:
\begin{codeexample}[code only]
\pgfkeys{/femm/measure/rotor torque/.style={quantity=force,method=band}}
\femmmeasure[rotor torque,name=at-60,parameter=60]
\end{codeexample}
|weighting| and |tolerance| configure the ordinary Laplace-weighted body method;
they do not change the band's prescribed weights. For an independent diagnostic,
register a closed contour strictly inside source-free air and use the existing
|method=maxwell,profile=...|. A single contour can have more radial bias than
the volume integral. Compare several contours and refined meshes.

Ordinary profiles and circuit measurements retain their syntax. For example,
|\femmplot[profile=gap-30,component=By,abscissa=angle]| inside a PGFPlots axis
plots the angular gap field. Set the profile centre explicitly for an offset
machine. Material selections use the original material names even when a
rotated anisotropic rotor internally needs its own tensor copy. Such a material
filter also selects stator regions with that material; add a geometric selection
when isolating a body for the ordinary weighted-force method.

\Needspace{10\baselineskip}
\subsection{Numerical construction and accuracy}
The template is meshed once with two protected polygonal traces. Band triangles
are discarded; retained bulk nodes and triangle connectivity are copied into
each position. Rotor coordinates, material axes and permanent-magnet directions
rotate physically. Uniform, radial, tangential and angular magnetisation laws
follow their material frame, including the law centre. Stator coordinates stay
unchanged. Each new annular strip is triangulated by connecting its ordered
inner and outer vertices. Orientation predicates, edge incidence, coverage,
Euler characteristic and retained constraints are checked before solving.

For |method=band|, the nodal weight $w$ is one on the rotor bulk, zero on the
stator bulk, and $1-\ell/L$ on intermediate ring $\ell$. P1 interpolation makes
its gradient constant in each band triangle. With planar Maxwell stress
\[
 \mathbf T=\frac1{\mu_0}\left(\mathbf B\mathbf B^{\mathsf T}
             -\frac12|\mathbf B|^2\mathbf I\right),\qquad
 \mathbf F=-\int_{\rm band}\mathbf T\nabla w\,\mathrm dV,
\]
the torque is the moment of the same force density about the chosen origin.
The piecewise-constant field and linear moment arm are integrated exactly for
the discrete P1 field. This exact integration does not eliminate finite-element
error. The band's geometric weights avoid an additional Laplace solve and
average stress over its thickness. They are not FEMM's AGE torque formula.

Useful checks include preservation of a uniform field for arbitrary angles,
zero net stress in that uniform field, a magnet in an imposed transverse field,
and the fixed-current identity
$T=\partial W'/\partial\theta$ with $\theta$ in radians. Compare the energy
finite difference at several angular steps; making the step extremely small
can amplify reconnection ripple. Refine bulk spacing, angular segmentation,
radial layers and exterior distance independently. The implementation does not
promise noise-free cogging torque or an error bound from a small solver residual.

\Needspace{10\baselineskip}
\subsection{Persistence, interchange and current limits}
Use separate cache files for the template mesh and each angle/current state.
The template can reuse a mesh-only cache; each position can reuse its solution.
Frozen mode checks the geometry, sources and exact inherited mesh just as it
does for other models. Native format 13 gains a moving-band signature carrying
method identity, template provenance, normalised angle, retained/reconnected
mesh and force weights. Ordinary models keep their previous fingerprints and
stored schema. Template names must be recreated in each TeX run; they are not
serialised as a document-wide object registry.

A supported position exports as an ordinary static |.fem| or |.ans| snapshot
at its current angle. Its material paths, magnet directions and mesh are already
rotated. A FEMM user can solve that one |.fem| normally; the file does not contain
a live moving-band controller or native AGE records. The answer export preserves
the computed P1 mesh and potentials. Existing representability checks still
apply: a rotated off-diagonal material tensor cannot be encoded in FEMM's
global diagonal material records, and nonlinear ANS export needs
|interpolation=femm|. Angular coercive laws retain their documented export rules.
Importing a snapshot creates an ordinary static model, not a sliding template.

The initial scope is one complete circular band in a rectangular native planar
DC problem with |mesher=delaunay| and |field model=open|. Custom domains,
tagged/paired boundaries, periodic sectors, imported templates, callbacks,
parallel/voltage circuits, harmonic excitation, ideal circuits and operating-point
responses are excluded. A native machine node may supply rotor/stator domains
when all its paths satisfy the air-collar contract. Existing TikZ node anchors
are drawing objects and are not recreated or rebound by a later position.

\Needspace{10\baselineskip}
\subsection{Lua entry points and validation}
\begin{codeexample}[code only]
local S = require("luafemm-sliding")
local template = S.capture(model, {
  inner_radius=4.3, outer_radius=5.3,
  center_x=0, center_y=0, segments=96, layers=3
})
local position = S.new(template, {
  angle=30, currents={drive=10},
  cache={file="machine-30",mode="auto"}
})
require("luafemm").solve(position)
local force = require("luafemm-post").measure(position, {
  name="at-30",quantity="force",method="band"
})
\end{codeexample}
|capture| meshes the supplied unmeshed model and returns independently owned
template data; treat those returned data as read-only. |new| returns an independent
meshed, unsolved model. Its angle defaults to zero, currents to an empty override
list and cache mode to |off|. Unknown options are errors. Physical edits to a
position after creation are rejected; create another position from the template.

The portable tests cover bulk identity, angle wrapping, affine fields, analytic
magnet torque, nonlinear coenergy differences, tensor and angular-law rotation,
source overrides and passive caches. The format matrix exercises shared cold,
warm and frozen data in plain LuaTeX, LuaLaTeX and ConTeXt MkIV. The optional
|scripts/sliding-interop.py| exports static angle snapshots and compares torque
and linkage with independent xfemm remeshes and solves. Numerical evidence and
its tolerances are recorded in \nolinkurl{docs/validation.md}.
